%--------------------------------------------------------------------------------
%
%
%--------------------------------------------------------------------------------
\pagebreak
\section*{EXTR Extract Bit Field}
\addcontentsline{toc}{section}{EXTR Extract Bit Field}

\instructionentry{Syntax}{
  \texttt{EXTR[.S] RegR, RegB, pos, len } \\
  \texttt{EXTR[.S] RegR, RegB, SAR, len }
}

\instructionentry{Format}{

	The EXTR instruction uses the immediate-20 instruction format. \\

   	\begin{flushleft}
    \begin{tikzpicture}[scale=0.7, transform shape]
        
        \draw[help lines, gray!70, dashed] (0,0) grid(16,1.5);
        
        \node[  tsRectangle, 
                minimum width=16cm,
                minimum height=1cm,
                text centered,
                fill=white!50] at (8,1) { };
                
     	\draw[-] (1,0.5) -- (1, 1.5);
     	\draw[-] (3,0.5) -- (3, 1.5);
     	\draw[-] (5,0.5) -- (5, 1.5);
     	
     	\draw[-] (6.5,0.5) -- (6.5, 1.5);
     	\draw[-] (8.5,0.5) -- (8.5, 1.5);
     	\draw[-] (9,0.5) -- (9,1.5);
     	\draw[-] (9.5,0.5) -- (9.5, 1.5);
     	\draw[-] (10,0.5) -- (10, 1.5);
     	\draw[-] (13,0.5) -- (13, 1.5);
     	
     	\node at (0.5,0.25) {2};
    	\node at (2,0.25) {4};
    	\node at (4,0.25) {4};
    	\node at (5.75,0.25) {3};
    	\node at (7.5,0.25) {4};
    	\node at (9,0.25) {2};
    	\node at (10.5,0.25) {4};
   	 	\node at (13.75,0.25) {9};
   	 	
   	 	\node at (0.5,1) {\small 0};
		\node at (2,1) {\small \OpCodeEXTR};
		\node at (4,1) {\small Reg R};
		\node at (5.75,1) {\small 0};
		
		\node at (7.5,1) {\small Reg B};
		\node at (8.75,1) {\small 0};
		\node at (9.25,1) {\small A};
		\node at (9.75,1) {\small S};
		\node at (11.5,1) {\small pos};
		\node at (14.5,1) {\small len };
   	 	
       \end{tikzpicture}
	\end{flushleft}
}

\instructionentry{Description}{

The \texttt{EXTR} instruction performs a bit field extract specified by the position and length instruction field from general register RegB. The position field specifies the rightmost bit of the bitfield to extract. The length field specifies the bit size if the field to extract. The extracted bit field is stored right justified in the general register RegR. If the S bit is set, the result is sign extended. If the A bit is set, the shift amount control register is used for obtaining the position value instead of the length instruction field.
}

\instructionentry{Operation}{

	\texttt{if ( instr.[A] tmp <- shamt; } \\
	\texttt{else tmp <- instr.[pos]; } \\
  	\texttt{RegR <- extractField( RegB, pos, len ); } \\
  	\texttt{if ( instr.[S] ) RegR <- signExtend( RegR, pos ); } \\
} 

\instructionentry{Exceptions}{

	None
}

\instructionentry{Notes}{

	None. 
}




